This study builds a geometry-constrained rock--TBM spatial interaction graph
from three inputs: a TSP-derived rock voxel field, an advancing parameterized
TBM surface, and excavation-step-aligned monitoring records. The core method
constructs this graph at each excavation step and aggregates weighted candidate
edges by TBM component into a component
$\times$ excavation-step anomalous interaction field. The resulting field
contains $I_c(t)$ and $\Delta I_c(t)$, can be aligned with monitoring
responses, and retains traceable edge contributions for excavation anomaly
indication, component-indexed readout, and jamming-risk-oriented
interpretation.

The two case studies indicate that the anomalous interaction field organises
rock--machine relations into interpretable component-indexed excavation-step
patterns and, most clearly in the SJLS case, retains residual-consistent
co-variation after simple trend-based null explanations are considered. The
component-level field also provides a step-change readout. In the SJLS case,
$\Delta I_c(t)$ has a large negative association with the step-level
shield-pressure residual change. Relation-support ablations indicate that,
within the present field-readout pairs, the empirical distinction is mainly
between component-indexed candidate-edge aggregation and pooled or global
summaries; the additional numerical effect of distance-normal weighting is
secondary. The evidence supports a spatial field-representation claim rather
than contact-force recovery, calibrated jamming
risk, unique component attribution, or broad predictive dominance.

The proposed field provides a spatial representation basis for future
rock--machine interaction analysis and digital-twin-oriented TBM excavation
analysis. In particular, the component $\times$ excavation-step anomalous
interaction field can be queried alongside monitoring responses to identify
which TBM component is exposed to anomalous rock conditions at which chainage.
It therefore offers a structured input for excavation anomaly indication and
jamming-risk-oriented interpretation in future excavation-analysis workflows.
